\title{Reproducible Research Through Explicit Design Models}
\subtitle{A worked thesis with bilingual abstracts and a complete academic apparatus}
\author{Alex Example}
\date{2026}
\SemioSetup{
  thesis-level=master,
  university={Example University},
  faculty={Faculty of Architecture},
  institute={Institute for Computational Design},
  degree={Master of Science},
  discipline={Computational Design},
  programme={Architecture and Technology},
  student-id={20260001},
  supervisor={Prof. Taylor Example},
  reviewers={Prof. Morgan Example},
  submission={2026-10-28},
  place={Hannover},
  keywords={reproducibility, design models, research methods}
}
\addbibresource{📚️references.bib}

\begin{document}
\makethesistitle
\begin{Abstract}
Reproducible research requires explicit assumptions, traceable sources and independently verifiable results. This worked thesis demonstrates a complete academic document, from its examination metadata to bilingual abstracts, numbered chapters, references and supplementary evidence. Its illustrative experiment compares two workflows under the same declared evaluation protocol.
\end{Abstract}
\begin{Kurzfassung}
Reproduzierbare Forschung benötigt explizite Annahmen, nachvollziehbare Quellen und unabhängig überprüfbare Ergebnisse. Diese Beispielarbeit zeigt den vollständigen Aufbau einer wissenschaftlichen Abschlussarbeit mit Prüfungsangaben, zweisprachigen Zusammenfassungen, Kapiteln, Quellen und ergänzenden Nachweisen.
\end{Kurzfassung}
\begin{Acknowledgements}
The author thanks the research participants and supervisors for their critical review of the study protocol. All names, institutions and numerical observations in this example are illustrative.
\end{Acknowledgements}
\makededication{For everyone who makes research understandable and reproducible.}
\begin{Abbreviations}
  \item[API] Application programming interface
  \item[CI] Continuous integration
  \item[PDF] Portable document format
\end{Abbreviations}
\tableofcontents
\makemainmatter

\chapter{Introduction}
\label{chap:introduction}
Research findings can only be checked when their assumptions and evaluation procedures are explicit. Literate programming offers one foundation for keeping explanations close to executable reasoning \cite{knuth1984}.
\section{Research Question}
How does an explicit model affect the repeatability of a design evaluation? The study measures the proportion of independently repeated evaluations that reproduce the declared result.
\section{Contributions}
The thesis contributes a documented protocol, an illustrative comparison and an account of the limits of that comparison. Supplementary material records the exact assumptions used in the experiment.

\chapter{Background and Related Work}
Scientific computing combines reproducible procedures with clear accounts of numerical and methodological limitations \cite{wilson2014}. This thesis treats a source, an evaluation protocol and its resulting observation as separate research artifacts.
\section{Definitions}
Repeatability means obtaining the same result under the same specified conditions. The present demonstration evaluates only those specified conditions; it does not establish general validity across institutions or datasets.

\chapter{Research Method}
\label{chap:method}
\section{Protocol}
Two workflows are evaluated on twenty illustrative tasks. Workflow A records informal instructions. Workflow B records the inputs, model version and evaluation steps explicitly. Every task is evaluated twice using the same protocol.
\section{Metric}
For $n$ evaluations, of which $r$ reproduce the declared result, the repeatability ratio is
\begin{equation}
  R = \frac{r}{n}, \qquad 0 \leq R \leq 1.
\end{equation}
\section{Ethics and Data Management}
The demonstration contains no personal participant data. A real study must replace these illustrative observations with its approved protocol, consent procedure and data availability statement.

\chapter{Results}
\begin{Table}[title=Illustrative repeatability observations]
\begin{tabular}{lrrr}
  Workflow & Reproduced & Evaluations & Ratio \\
  A & 14 & 20 & 0.70 \\
  B & 19 & 20 & 0.95
\end{tabular}
\end{Table}
The absolute difference in this illustrative dataset is $0.25$. The observations demonstrate the reporting format rather than establish an empirical claim.
\begin{Figure}[title=Research evidence chain,height=3cm,anchor=center,justify=center]
  Declared inputs $\longrightarrow$ explicit model $\longrightarrow$ evaluation protocol $\longrightarrow$ archived result
\end{Figure}

\chapter{Discussion}
The protocol in Chapter~\ref{chap:method} makes the comparison inspectable. The small constructed dataset cannot support generalization or causal inference. A subsequent study should preregister its hypotheses, report uncertainty and compare results across independent teams.
\section{Threats to Validity}
Task selection, evaluator experience and tooling differences may affect repeatability. These factors should be controlled or reported rather than inferred from the aggregate ratio.

\chapter{Conclusion and Future Work}
The research question in Chapter~\ref{chap:introduction} motivates a transparent evaluation protocol. The worked document connects its question, method, observations and limitations, and archives the references and supplementary material needed for review.

\makeappendices
\appendixchapter[PR]{Evaluation Protocol}
Record each task identifier, input version, evaluator, expected result and observed result. Preserve failed evaluations alongside successful ones. Compute the ratio using all twenty tasks, without excluding inconvenient observations.
\appendixchapter[DA]{Data Availability}
All numerical values in this example are authored demonstration data. Replace this statement with the persistent location, access conditions and license of the actual research materials.
\appendixchapter{References and Registers}
\makeregisters
\makedeclaration
\end{document}
